\documentclass[12pt,a4paper]{article}
\usepackage[utf8]{inputenc}
\usepackage[portuguese]{babel}
\usepackage{amsmath}
\usepackage{geometry}
\usepackage{listings}
\usepackage{xcolor}
\usepackage{hyperref}

\geometry{margin=2.5cm}

\lstset{
    language=Python,
    basicstyle=\ttfamily\small,
    keywordstyle=\color{blue}\bfseries,
    stringstyle=\color{orange},
    commentstyle=\color{gray}\itshape,
    backgroundcolor=\color{gray!10},
    frame=single,
    breaklines=true,
    showstringspaces=false,
    numbers=left,
    numberstyle=\tiny\color{gray},
    numbersep=8pt
}

\title{
    \Large\textbf{Universidade de Fortaleza - UNIFOR}\\
    \large Curso de Ciência da Computação\\
    \vspace{1.5cm}
    \huge\textbf{Documentação da Implementação}\\
    \large Análise de Rede de Colaboração Científica (ca-GrQc)
}
\author{João Pedro Brasil, Lucas Lustosa Dias, Davi Dias Vale}
\date{12 de fevereiro de 2026}

\begin{document}

\maketitle
\thispagestyle{empty}
\newpage

\tableofcontents
\newpage

\section{Visão Geral}

Este documento descreve a implementação da classe \texttt{Graph} contida em \texttt{graph.py}, responsável por representar e analisar o grafo de colaboração científica ca-GrQc. A implementação é independente de bibliotecas de grafos externas para as operações estruturais, utilizando apenas estruturas de dados fundamentais.

\section{Estrutura de Dados}

\subsection{Representação do Grafo}

O grafo é representado internamente por uma \textbf{lista de adjacência}. Cada vértice possui um \texttt{Bag} (multiconjunto ordenado da biblioteca \texttt{algs4}) que armazena seus vizinhos.

\begin{lstlisting}
def __init__(self, v):
    self.V = v      # numero de vertices
    self.E = 0      # numero de arestas
    self.map = {}   # ID original -> posicao interna
    self.adj = []   # lista de adjacencia (lista de Bags)
\end{lstlisting}

\textbf{Por que lista de adjacência e não matriz?}

Com $|V| = 5.242$ vértices, uma matriz de adjacência teria $5.242^2 \approx 27,5$ milhões de células — a grande maioria zerada, já que a rede é esparsa (densidade de 0,21\%). A lista de adjacência usa memória proporcional apenas às arestas existentes: $O(|V| + |E|)$.

\subsection{Mapeamento de Vértices}

Os IDs dos pesquisadores no arquivo de dados são inteiros arbitrários (ex: 3466, 17038). O dicionário \texttt{map} traduz cada ID original para uma posição interna sequencial na lista \texttt{adj}:

\begin{lstlisting}
def add_mapping(self, node_value: int) -> None:
    self.map[node_value] = len(self.adj)
    self.adj.append(Bag())

def get_position(self, node_value: int) -> int:
    return self.map[node_value]

def get_value(self, node_position: int) -> int:
    for key, value in self.map.items():
        if value == node_position:
            return key
\end{lstlisting}

Isso permite preservar os IDs originais dos pesquisadores sem exigir que sejam sequenciais ou contíguos.

\section{Inserção de Arestas}

\subsection{Leitura do Arquivo}

O método \texttt{add\_node\_from\_file} lê o arquivo de dados linha a linha, tratando separadores de tabulação e quebras de linha, e chama \texttt{add\_edge} para cada par de vértices:

\begin{lstlisting}
def add_node_from_file(self, file_path: str = "data.txt") -> None:
    with open(file_path, 'r') as file:
        linhas = file.readlines()
    for linha in linhas:
        linha_tratada = linha.replace("\t"," ").replace("\n","").split(" ")
        no_origem, no_destino = map(int, linha_tratada)
        self.add_edge(no_origem, no_destino)
\end{lstlisting}

\subsection{Inserção Bidirecional}

Como o grafo é \textbf{não-direcionado}, cada aresta é inserida nos dois sentidos. Vértices são criados dinamicamente na primeira vez que aparecem:

\begin{lstlisting}
def add_edge(self, v, w, directed: bool = False):
    if v not in self.map:
        self.add_mapping(v)
    if w not in self.map:
        self.add_mapping(w)
    if directed:
        if w in self.adj[self.get_position(v)]:
            return
    v_position = self.get_position(v)
    w_position = self.get_position(w)
    self.adj[v_position].add(w)  # v -> w
    self.adj[w_position].add(v)  # w -> v
    self.E += 1
\end{lstlisting}

O parâmetro \texttt{directed} permite uso futuro do grafo de forma direcionada. No caso do ca-GrQc, é sempre \texttt{False}.

\section{Cálculo de Métricas}

\subsection{Grau de um Vértice}

O grau é obtido diretamente pelo tamanho do \texttt{Bag} de adjacência do vértice:

\begin{lstlisting}
def degree(self, v, from_value: bool = True) -> int:
    v_position = self.get_position(v) if from_value else v
    return self.adj[v_position].size()
\end{lstlisting}

O parâmetro \texttt{from\_value} permite consultar pelo ID original (\texttt{True}) ou pela posição interna (\texttt{False}).

\subsection{Densidade}

Implementa diretamente a fórmula $d = \frac{2|E|}{|V|(|V|-1)}$:

\begin{lstlisting}
def density(self):
    return (2 * self.E) / (self.V * (self.V - 1))
\end{lstlisting}

O denominador $|V|(|V|-1)$ é o número máximo de arestas em um grafo simples não-direcionado com $|V|$ vértices.

\subsection{Grau Mínimo e Máximo}

Calculados iterando sobre todos os vértices do mapeamento:

\begin{lstlisting}
def min_degree(self) -> int:
    return min(self.degree(v) for v in self.map.keys())

def max_degree(self) -> int:
    return max(self.degree(v) for v in self.map.keys())
\end{lstlisting}

\subsection{Distribuição de Graus}

Retorna um dicionário \texttt{\{grau: frequência\}} usando \texttt{Counter}:

\begin{lstlisting}
def get_degree_distribution(self):
    degrees = [self.degree(v) for v in self.map.keys()]
    return dict(Counter(degrees))
\end{lstlisting}

\section{Ajuste de Lei de Potência}

\subsection{Fundamentação Matemática}

Uma rede scale-free tem distribuição de graus $P(k) \sim k^{-\gamma}$. Aplicando logaritmo:

\begin{equation}
    \log P(k) = \log C - \gamma \cdot \log k
\end{equation}

Isso transforma o problema em uma regressão linear no espaço log-log, onde $\gamma$ é o módulo do coeficiente angular da reta ajustada.

\subsection{Implementação}

O método \texttt{fit\_power\_law} executa as seguintes etapas:

\textbf{1. Normalização para probabilidade:}
\begin{lstlisting}
degree_dist = self.get_degree_distribution()
degrees = np.array(sorted(degree_dist.keys()))
frequencies = np.array([degree_dist[k] for k in degrees])
probabilities = frequencies / frequencies.sum()
\end{lstlisting}

\textbf{2. Filtragem de zeros (evita indefinição no log):}
\begin{lstlisting}
mask = (probabilities > 0) & (degrees > 0)
degrees_filtered = degrees[mask]
probabilities_filtered = probabilities[mask]
\end{lstlisting}

\textbf{3. Regressão linear no espaço log-log:}
\begin{lstlisting}
log_k = np.log(degrees_filtered)
log_p = np.log(probabilities_filtered)

coeffs = np.polyfit(log_k, log_p, 1)
gamma = -coeffs[0]  # coeficiente angular negativo = gamma
log_C = coeffs[1]
C = np.exp(log_C)
\end{lstlisting}

\textbf{4. Cálculo do $R^2$:}
\begin{lstlisting}
log_p_pred = log_C - gamma * log_k
ss_res = np.sum((log_p - log_p_pred)**2)
ss_tot = np.sum((log_p - log_p.mean())**2)
r_squared = 1 - (ss_res / ss_tot)
\end{lstlisting}

O $R^2$ mede qual fração da variância de $\log P(k)$ é explicada pelo modelo linear. Valores acima de 0,8 indicam indício razoável de lei de potência; acima de 0,9, indício forte.

\subsection{Interpretação do Expoente}

O método \texttt{interpret\_gamma} classifica o valor de $\gamma$ conforme os intervalos estabelecidos na literatura:

\begin{lstlisting}
def interpret_gamma(self, gamma: float) -> str:
    if gamma < 2:
        return "gamma < 2: Rede extremamente concentrada em hubs."
    elif 2 <= gamma < 3:
        return "2 <= gamma < 3: Rede scale-free tipica (media finita, variancia infinita)."
    elif 3 <= gamma < 4:
        return "3 <= gamma < 4: Rede com hubs menos dominantes, mas ainda presentes."
    else:
        return "gamma >= 4: Rede sem hubs bem definidos."
\end{lstlisting}

\textbf{Nota:} a primeira condição (\texttt{gamma < 2}) contém um erro tipográfico na mensagem original — exibe ``$\gamma > 2$'' em vez de ``$\gamma < 2$''. Isso não afeta os resultados, pois o $\gamma$ obtido (2,04) nunca aciona esse ramo.

\section{Visualização}

\subsection{Histograma de Graus}

Gera e salva o histograma da distribuição de graus:

\begin{lstlisting}
def save_histogram_degree(self, save_path: str = "histogram.png") -> None:
    degrees = [self.degree(v) for v in self.map.keys()]
    plt.hist(degrees, bins=range(max(degrees) + 1),
             color='lightblue', edgecolor='black')
    plt.title('Histograma de Graus dos Vertices')
    plt.xlabel('Grau')
    plt.ylabel('Frequencia')
    plt.savefig(save_path)
\end{lstlisting}

\subsection{Gráfico Log-Log com Ajuste}

Gera o gráfico log-log com os dados observados e a curva ajustada, anotando a interpretação de $\gamma$ diretamente no gráfico:

\begin{lstlisting}
def save_power_law_plot(self, save_path: str = "power_law_fit.png"):
    fit_results = self.fit_power_law()
    gamma = fit_results['gamma']
    C = fit_results['C']
    r_squared = fit_results['r_squared']
    degrees = fit_results['degrees']
    probabilities = fit_results['probabilities']

    plt.loglog(degrees, probabilities, 'bo', alpha=0.6,
               label='Dados observados')

    k_fit = np.linspace(degrees.min(), degrees.max(), 100)
    p_fit = C * k_fit**(-gamma)
    plt.loglog(k_fit, p_fit, 'r-', linewidth=2,
               label=f'P(k) ~ k^(-gamma), gamma={gamma:.3f}, R2={r_squared:.4f}')

    interpretation = self.interpret_gamma(gamma)
    plt.text(0.02, 0.02, interpretation, transform=plt.gca().transAxes,
             fontsize=9, bbox=dict(boxstyle='round', facecolor='wheat'))

    plt.savefig(save_path, dpi=300, bbox_inches='tight')
\end{lstlisting}

\end{document}
